\section*{Chapitre \ref{ch:nw:where-the-crypto-venues-live} --- Où vivent les plateformes crypto}

\begin{solution}{exo:nw:where-the-crypto-venues-live:1}
Au plus $0.5~\mathrm{ms} \times c_0 / 1.462 = \qty{102.5}{\kilo\metre}$ le long de la géodésique, et moins une fois le matériel et le traitement retranchés.
\end{solution}

\begin{solution}{exo:nw:where-the-crypto-venues-live:2}
Parce que l'adresse peut appartenir au \omterm{def:nw:where-the-crypto-venues-live:as}{réseau de périphérie} (edge) du fournisseur (atteint par \omterm{def:nw:where-the-crypto-venues-live:as}{anycast}), qui termine les connexions près de l'utilisateur et les
transmet à l'origine~: la recherche indique alors «~edge global~», et non la région du moteur d'appariement.
\end{solution}

\begin{solution}{exo:nw:where-the-crypto-venues-live:3}
Le meilleur temps attendu passe de 279 à \qty{185}{\micro\second} sur les dix premiers lancements, soit \qty{94}{\micro\second}~; de quarante à 160 lancements, il passe
de 157 à 138, soit \qty{19}{\micro\second}.
\end{solution}

\begin{solution}{exo:nw:where-the-crypto-venues-live:4}
Les zones sont distantes de quelques kilomètres, soit quelques microsecondes de lumière, alors que l'incertitude d'une triangulation depuis des milliers de
kilomètres (\omterm{def:nw:the-north-american-map:factor}{facteurs de tracé}, matériel, bruit) se chiffre en dizaines ou en centaines de kilomètres. À l'intérieur de la région, seules les machines
de chaque zone qui mesurent la plateforme directement peuvent distinguer les zones.
\end{solution}

\begin{solution}{exo:nw:where-the-crypto-venues-live:5}
Son aller-retour vers la plateforme est passé d'une région à l'autre, à plus de \qty{200}{\milli\second} à travers le monde~: elle ne pouvait plus coter de
façon compétitive. Elle aurait dû suivre la plateforme~: lire l'annonce, lancer des machines dans la nouvelle région avant la bascule, y faire tourner la
loterie et déplacer son moteur de cotation.
\end{solution}

\begin{solution}{exo:nw:where-the-crypto-venues-live:6}
Il accumule des indices jusqu'à ce que la trajectoire franchisse une frontière calibrée pour un taux de fausses alarmes de 5~\% sur l'ensemble de
l'échantillon~; il faut donc attendre de nombreux jours d'indices, ici douze jours pour un décalage de quatre écarts types. Il faut le combiner avec un
seuil rapide sur la médiane quotidienne ou horaire, confirmé par une seconde mesure, et garder le test CUSUM pour confirmer et pour détecter les dérives lentes.
\end{solution}

\begin{solution}{exo:nw:where-the-crypto-venues-live:7}
\texttt{best\_n(LOTTERY, 20.0, 17.14)} donne 23 lancements, valant environ 1\,830~USD par mois nets~: l'optimum se déplace vers le bas et le gain est divisé
par six.
\end{solution}

\begin{solution}{exo:nw:where-the-crypto-venues-live:8}
L'API peut être servie par un \omterm{def:nw:where-the-crypto-venues-live:as}{réseau de périphérie} qui répond depuis un \omterm{def:nw:the-european-map:pop}{point de présence} proche de Francfort et transmet à l'origine ailleurs~; le
temps d'aller-retour mesure la distance jusqu'à la périphérie, et non jusqu'au moteur d'appariement. Il faut mesurer le trajet réellement parcouru par les
ordres (le chemin de saisie d'ordres, ou le temps écoulé entre un ordre et son accusé de réception), et comparer l'adresse aux plages du fournisseur.
\end{solution}

\begin{solution}{pb:nw:where-the-crypto-venues-live:1}
\textbf{Partie I.}
\begin{enumerate}
\item \qty{279}{\micro\second} (la moyenne d'une loi log-normale dépasse sa médiane).
\item 185, 157 et \qty{138}{\micro\second}.
\item La loi log-normale est asymétrique à droite~: les lancements lents et rares tirent la moyenne au-dessus de la médiane.
\item Chaque lancement supplémentaire fait gagner moins que le précédent~: le meilleur temps attendu baisse de plus en plus lentement, si bien que l'optimum est fini.
\end{enumerate}
\textbf{Partie II.}
\begin{enumerate}[resume]
\item Environ 11\,400~USD par mois.
\item 77 lancements, soit environ 11\,700~USD par mois nets.
\item Très plat~: 40 et 120 lancements sont à quelques~\% près de l'optimum.
\item 23 lancements, soit environ 1\,830~USD par mois.
\end{enumerate}
\textbf{Partie III.}
\begin{enumerate}[resume]
\item Faits~: la médiane publiée de l'aller-retour à l'intérieur d'une zone (\qty{270}{\micro\second}, chapitre~16) et le chiffre «~sous 150~µs~» du blog d'AWS pour l'intra-région.
  Hypothèses~: qu'il s'agit de la médiane et du meilleur 1~\% des lancements, la forme log-normale, l'indépendance, la valeur et le coût.
\item Le fournisseur peut placer des lancements successifs près les uns des autres (ou volontairement à l'écart), et la capacité varie au cours de la journée.
\item Le type d'instance et sa \omterm{def:nw:networking-for-trading:frame}{carte réseau}, son réglage et sa pile de contournement du noyau (kernel bypass), les autres locataires de l'hôte, ainsi que les serveurs et les
  répartiteurs de charge de la plateforme.
\item Lancer les machines, mesurer la distribution de l'aller-retour de chacune avec des horodatages matériels, ajuster la distribution, et recommencer à des jours différents.
\end{enumerate}
\textbf{Partie IV.}
\begin{enumerate}[resume]
\item \emph{Résultat nommé}~: avec la distribution du modèle, le meilleur aller-retour attendu passe de \qty{279}{\micro\second} pour un lancement à 185 pour dix,
  157 pour quarante et 138 pour 160~; à 100~USD par mois et par microseconde et 17.14~USD par lancement, 77 lancements maximisent la valeur nette, soit environ 11\,700~USD par mois.
\item Tant que la plateforme ou le fournisseur ne change rien~: il faut mesurer à nouveau chaque jour et relancer la loterie sur un signal.
\item La zone~: la loterie ne tourne alors que dans la bonne zone, et personne ne gaspille de lancements dans les autres.
\item Il est ouvert à tout utilisateur qui paie ses lancements, mais il avantage ceux qui peuvent en financer beaucoup~; une plateforme qui veut un accès équitable
  l'égaliserait, comme le font les plateformes du chapitre~9.
\item Proposer des points d'accès à latence égale (des points de terminaison privés dans chaque zone, aux chemins égaux, chapitre~18), ou ajouter un délai qui supprime la différence.
\item Rien~: il lance une machine dans la bonne zone et accepte la médiane.
\item Dans les lignes du cloud~: la zone, les machines conservées, la routine de nouvelle mesure et le budget de lancements.
\item Parce que le fournisseur place les machines à des distances aléatoires de la plateforme, et que le meilleur parmi quarante vaut plus que le coût des quarante lancements.
\end{enumerate}
\end{solution}

\begin{solution}{iq:nw:where-the-crypto-venues-live:1}
Lire les déclarations publiques de la plateforme et du fournisseur~; comparer ses adresses de point de terminaison aux plages publiées par le fournisseur
(en distinguant le \omterm{def:nw:where-the-crypto-venues-live:as}{réseau de périphérie} des origines)~; trianguler depuis des machines situées dans plusieurs régions~; et confirmer en lançant des machines dans la région candidate et en mesurant.
\iqlookfor{Sources datées~; plages IP et périphéries~; la mesure comme dernier mot.}
\end{solution}

\begin{solution}{iq:nw:where-the-crypto-venues-live:2}
Une seule adresse annoncée depuis de nombreux endroits, si bien que chaque paquet va vers le site le plus proche dans le système de routage. Un aller-retour vers une adresse \omterm{def:nw:where-the-crypto-venues-live:as}{anycast} mesure la
distance jusqu'à l'instance la plus proche, qui peut être une périphérie transmettant à une origine lointaine.
\iqlookfor{Le routage choisit le site~; périphérie contre origine~; mesurer le trajet que suivent les ordres.}
\end{solution}

\begin{solution}{iq:nw:where-the-crypto-venues-live:3}
Garder autant de machines que la stratégie en a besoin (avec une de réserve), en choisissant la meilleure selon un centile de queue et non un simple ping~; fixer le nombre de lancements en
comparant le gain attendu d'un lancement supplémentaire à son coût, et re-mesurer chaque jour les machines conservées.
\iqlookfor{Statistiques d'ordre~; la queue plutôt que la moyenne~; coût de la recherche~; nouvelle mesure.}
\end{solution}

\begin{solution}{iq:nw:where-the-crypto-venues-live:4}
Suivre la médiane quotidienne (ou horaire) de l'aller-retour de chaque machine vers la plateforme~; signaler un saut au-delà d'un seuil et le confirmer~; lancer un test CUSUM pour les dérives lentes~;
surveiller aussi les annonces de la plateforme et les adresses vers lesquelles ses points de terminaison se résolvent.
\iqlookfor{Alarme rapide plus confirmation statistique~; annonces~; adresses.}
\end{solution}

\begin{solution}{iq:nw:where-the-crypto-venues-live:5}
Les heures de négociation asiatiques concentrent une part importante du volume crypto, plusieurs grandes plateformes tournent dans la région Tokyo d'AWS, et des sociétés de trading et d'autres plateformes ont suivi~: un
effet de réseau, que BitMEX a invoqué lorsqu'il a quitté l'Irlande, et qui donne aux traders proches un avantage de latence de plusieurs centaines de millisecondes sur les
traders éloignés.
\iqlookfor{Effet de réseau~; exemples sourcés~; asymétrie de latence entre régions.}
\end{solution}

\begin{solution}{iq:nw:where-the-crypto-venues-live:6}
Chaque aller-retour borne la distance (la moitié du trajet à la vitesse de la lumière dans le milieu) et, avec un \omterm{def:nw:the-north-american-map:factor}{facteur de tracé}, l'estime~; trouver le point qui s'accorde le mieux
avec les estimations à l'intérieur de chaque borne. La précision est limitée par les \omterm{def:nw:the-north-american-map:factor}{facteurs de tracé} inconnus et les délais matériels, par le bruit, et par la
géométrie des sondes~: à des milliers de kilomètres, elle se chiffre en dizaines ou centaines de kilomètres.
\iqlookfor{Bornes contre estimations~; sensibilité au \omterm{def:nw:the-north-american-map:factor}{facteur de tracé}~; géométrie des sondes.}
\end{solution}
